\begin{abstract}
Pretrained perception models---depth estimators, object detectors, trackers, and vision-language models---are now standard components in robot learning systems, yet the choice of which model to deploy is made without evidence from the conditions robots actually face.
Standard benchmarks use curated internet images; robots see the world through low-resolution Intel RealSense sensors under clutter, hand occlusion, motion blur, and challenging surfaces.
We introduce \coolname{}, a real-world evaluation methodology and benchmark that quantifies the \textit{deployment readiness} of pretrained perception models for embodied robotics.
%
A human carrying the same sensor rig a mobile robot would use (D435 RGB-D + T265 tracking + GoPro egocentric) captures 100 scenes spanning indoor and outdoor environments, daylight and night, reflective and transparent surfaces, across ${\sim}$70 object categories and 75{,}000 frames.
%
Each scene is captured in three phases that mirror the robot manipulation cycle:
\textbf{clutter} (the scene as a robot naturally finds it),
\textbf{interaction} (a human manipulates objects---the data that learning-from-humans pipelines train on), and
\textbf{clean} (the scene reorganised---the condition models are designed for).
%
The gap between clean-scene and clutter/interaction performance \textit{is} the deployment readiness gap, and no existing benchmark measures it.
%
\coolname{} evaluates \todo{N} models across \todo{M} perception tasks on the \textit{same} physical scenes and objects, using deployment-readiness metrics---State-Transition Robustness, Geometric Coherence, and Effort-Stratified Difficulty---that capture properties invisible to static accuracy and compose into an Embodied Deployment Score (EDS).
%
\todo{HEADLINE FINDING.}
%
\coolname{} is released as a pip-installable, bring-your-own-model toolkit at \url{\webpage}.
\end{abstract}
